\section{Relation to prior work and remaining limits}
\label{sec:context}

\subsection{Exact wide-network descriptions carry history}

Tensor Programs rigorously identify discrete-time infinite-width
feature-learning limits at each fixed finite training horizon
\citep{yang2021tensor}.  In a generic deep nonlinear network, the exact
construction introduces Gaussian variables, covariances, and transpose
response terms indexed by previous steps.  Yang and Hu describe a generic
quadratic-in-step storage and computation burden for this representation.
Recent work proves rich feature nondegeneracy throughout deep $\mu$P training,
while its recursive state continues to contain historical dependencies
\citep{chen2025global}.  These results establish that the standard exact
representation grows with its horizon; they do not prove that every exact
realization must have growing memory.

Self-consistent neural-network DMFT organizes infinite-width feature learning
through two-time activation, gradient, and causal response kernels
\citep{bordelon2022selfconsistent}.  There are also rigorous DMFT convergence
theorems for random-design high-dimensional first-order methods
\citep{celentano2026highdimensional,gerbelot2024rigorous}.  Those theorems take
sample and ambient dimensions to infinity proportionally while keeping a
finite internal learner dimension; they do not directly identify the
fixed-data, all-hidden-width deep regime here.  It would therefore be wrong to
describe DMFT as uniformly nonrigorous, just as it would be wrong to import a
random-design convergence theorem into this network without checking its
regime.

Rigorous multilayer mean-field limits and finite-width couplings are known
under mean-field parametrizations and paper-specific initialization structures
\citep{nguyen2023rigorous,araujo2019meanfield}.  Earlier measure-gradient-flow
work gives foundational shallow particle limits \citep{chizat2018global,
sirignano2020meanfield}.  Our theorem addresses a different question: it
compresses the learned dense increment along a fixed realized nonlinear
gradient-flow trajectory, then controls the induced feedback error.

\subsection{Polynomial memory and low-rank training}

HiPPO is direct prior art for the central memory mechanism
\citep{gu2020hippo}.  It derives online ODEs for optimal polynomial projections
of a signal history, including expanding-interval Legendre coordinates and
approximation bounds.  We do not claim to introduce online Legendre memory.
The additional ingredients here are the identification of paired
forward/backward neural histories, their reconstruction of every internal
learned matrix, an autonomous self-consistent training loop, and a
finite-horizon stability proof for the perturbed trajectory.

Retaining $W_0$ and representing $W_t-W_0$ in low rank also has precedent.
InRank studies cumulative low-rank updates and proves modewise behavior in a
restricted deep-linear setting \citep{zhao2024inrank}; invariant-subspace
compressions have rigorous deep-linear results
\citep{kwon2024compression,yaras2024compressible}.  The response factors here
are neither assumed to lie in a fixed subspace nor optimized as free low-rank
parameters.  They are causal summaries of the dense update's actual response
sources, and their directions evolve as the network learns.

Among the primary sources reviewed for this draft, we did not find a theorem
with the following conjunction: a fixed-order Legendre closure of coupled
forward/backward response histories for deep nonlinear neural-network gradient
flow, together with a finite-horizon trajectory error bound for the closed
system at fixed finite width.  This is a search-qualified distinction, not an
unconditional priority claim.

\subsection{What the result does and does not compress}

The strongest exact statement is deliberately narrow:
\begin{itemize}[leftmargin=*]
\item $P$ is fixed throughout one closure run, so state does not grow with
      elapsed time.
\item The learned dense correction is replaced by population moment fields and
      has a finite-rank action.
\item The initialized hidden matrices are retained as fixed operators, in both
      forward and transpose orientations.
\item Approximation is global on every prescribed finite horizon after choosing
      sufficiently large $P$; it is not uniform over all $t\ge0$.
\item The theorem is fixed-width.  Corollary~\ref{cor:diagonal} gives a
      width-dependent diagonal bridge, not a fixed-order population theorem.
\end{itemize}

These distinctions explain why the result can be both a real training
compression and an incomplete model compression.  It removes the dense
\emph{moving} internal matrices and unbounded historical state.  It does not
remove the random operator that repeatedly mixes neuron populations.  A future
width-uniform theory must either characterize that reused operator by a causal
population law or prove an effective-rank approximation compatible with both
orientations.

The sample index is another unresolved axis.  Current moving state grows as
$mP$ per layer.  Arbitrary labels preclude a universally sample-independent
exact state even in generic fixed-kernel dynamics.  Structured input laws may
permit a second approximation, such as Fourier--Galerkin compression of the
sample-indexed population fields, but that is separate from the theorem here.

\subsection{Open mathematical and empirical questions}

Several questions remain concrete.

\begin{enumerate}[leftmargin=*]
\item \textbf{Width-uniform estimates.}  The constants in
      Theorem~\ref{thm:main} may depend strongly on $n$.  The Euclidean
      response-arclength clock also scales with the concatenated neuron state.
      A population-normalized clock and width-uniform stability estimate are
      needed for a direct $n\to\infty$ theorem at controlled order.
\item \textbf{All-time stability.}  The current bounds contain finite-horizon
      clock lengths and a Gr\"onwall factor.  A uniform all-time result would
      need finite total activity or a trained propagator whose integrated
      amplification remains bounded.
\item \textbf{Initialized-operator compression.}  Compressing $W_0$ while
      preserving its correlated reuse and its adjoint is essential for a true
      total-memory improvement.  Fresh random redraws do not preserve the
      target dynamics.
\item \textbf{Nonsmooth and normalized networks.}  Exact ReLU, normalization
      layers, biases, convolution, residual connections, stochastic gradients,
      and discrete optimizers require new analysis.  Existing numerical tests
      outside the theorem show both successes and failures.
\item \textbf{Practical clock design.}  The response clock proves a better
      asymptotic order but introduces a Gram system and can behave worse at
      small $P$.  A clock optimized for projection error, numerical
      conditioning, and population scaling remains open.
\item \textbf{Statistical validation.}  The present experiments mostly use one
      initialization and matched-loss endpoints.  Seed distributions,
      continuous-time solver certificates, width trends, and broader tasks are
      needed before making general performance claims.
\end{enumerate}

At a conceptual level, a deterministic population learning law could play a
role analogous to an idealized model of computation: finite networks would be
understood through the ideal law plus a quantitative finite-resource error.
That analogy is a research program rather than a conclusion of this paper.
The response-memory theorem supplies one missing component---a controlled
finite description of learned history at fixed width---while leaving the
operator and population-limit steps visible.

